\chapter[Metodologia]{Metodologia}
\label{cap:metodologia}

Este capítulo descreve como a pesquisa é conduzida. A
Seção~\ref{sec:classificacao} classifica a pesquisa. A
Seção~\ref{sec:dsr} apresenta a \textit{Design Science Research} e o
mapeamento de suas atividades para este trabalho. A
Seção~\ref{sec:contexto} descreve o contexto em que o artefato é
desenvolvido. A Seção~\ref{sec:avaliacao} define os critérios e
procedimentos de avaliação, e a Seção~\ref{sec:ferramentas} lista as
ferramentas utilizadas.

\section{Classificação da pesquisa}
\label{sec:classificacao}

Quanto à natureza, a pesquisa é \textbf{aplicada}: busca resolver um problema
concreto de uma plataforma de dados governamentais em operação. Quanto à
abordagem, é \textbf{mista}: combina medidas quantitativas (métricas do
modelo, reprodutibilidade de execuções) com análise qualitativa da aderência
da arquitetura a princípios da literatura. Quanto aos procedimentos, segue a
\textbf{Design Science Research}, apoiada em pesquisa bibliográfica e em
pesquisa documental sobre o sistema existente.

\section{Design Science Research}
\label{sec:dsr}

A DSR é um paradigma de pesquisa em que o conhecimento é produzido pela
construção e avaliação de artefatos que resolvem problemas relevantes
\cite{hevner2004}. Artefatos podem ser construtos, modelos, métodos ou
instanciações. Neste trabalho, o artefato principal é um \textbf{modelo} (a
arquitetura de MLOps), acompanhado de uma \textbf{instanciação} (sua
implementação no caso de prevenção de paralisação de obras).

\citeonline{hevner2007} descreve a DSR em três ciclos. O \textbf{ciclo de
relevância} conecta o artefato ao ambiente, de onde vêm requisitos e para onde
voltam os resultados. O \textbf{ciclo de rigor} conecta o artefato à base de
conhecimento: teorias, métodos e trabalhos anteriores. O \textbf{ciclo de
design} itera entre a construção e a avaliação do artefato. Aqui, o ambiente
é o Gov Hub Cidades e o MCID, a base de conhecimento é a literatura do
Capítulo~\ref{cap:referencial} e o ciclo de design corresponde às iterações
de desenvolvimento da arquitetura.

Para operacionalizar o processo, adota-se a metodologia de
\citeonline{peffers2007}, que organiza a DSR em seis atividades. A
Tabela~\ref{tab:dsr} mostra como cada atividade se realiza neste trabalho e em
qual etapa (TCC1 ou TCC2) ela ocorre.

\begin{table}[htb]
	\centering
	\caption{Atividades da DSR \cite{peffers2007} aplicadas a este trabalho.}
	\label{tab:dsr}
	\small
	\begin{tabularx}{\textwidth}{p{3.4cm}Xc}
		\toprule
		\textbf{Atividade} & \textbf{Realização neste trabalho} & \textbf{Etapa} \\
		\midrule
		1. Identificar o problema e motivar &
		Diagnóstico da regra de gargalo vigente, levantamento dos dados
		disponíveis e da ausência de infraestrutura de ML na plataforma
		(Capítulos~\ref{cap:introducao} e~\ref{cap:proposta}) & TCC1 \\
		2. Definir os objetivos da solução &
		Requisitos da arquitetura derivados da literatura e das regras da
		plataforma (Seção~\ref{sec:requisitos}) & TCC1 \\
		3. Projetar e desenvolver &
		Projeto da arquitetura (ADR) e implementação dos componentes:
		transformações de painel, rótulo e features; código de treino,
		avaliação e escore; fluxos de orquestração & TCC1/TCC2 \\
		4. Demonstrar &
		Execução do ciclo completo no caso de prevenção de paralisação de obras
		do MCMV & TCC2 \\
		5. Avaliar &
		Critérios AV1--AV5 (Seção~\ref{sec:avaliacao}) & TCC2 \\
		6. Comunicar &
		Este documento, a monografia final e os ADRs e a documentação no
		repositório do projeto & TCC1/TCC2 \\
		\bottomrule
	\end{tabularx}
	\fonte{Elaborado pelos autores.}
\end{table}

O ponto de entrada da pesquisa é centrado no problema
\cite{peffers2007}: parte-se de uma necessidade observada no ambiente (antecipar
paralisações) e de uma lacuna de engenharia (ausência de um ciclo de vida de
ML integrado à plataforma).

\section{Contexto da pesquisa}
\label{sec:contexto}

O trabalho é desenvolvido no LabLivre da Universidade de Brasília,
no âmbito do projeto Gov Hub, que constrói plataformas de dados para órgãos do
governo federal. O ambiente de demonstração é a plataforma de dados do
Ministério das Cidades, organizada como um \textit{Data Lakehouse} em camadas
medalhão (Figura~\ref{fig:govhub}), que contém:

\begin{itemize}
	\item fluxos orquestrados de ingestão de fontes como o SFTP da Caixa, as
	planilhas do Novo MCMV, o histórico da Secretaria Nacional de Habitação
	(SNH) e APIs públicas, como as do IBGE e do Banco Central;
	\item uma zona de pouso em armazenamento de objetos, dividida em
	\texttt{raw/}, com os arquivos no formato original da fonte, e
	\texttt{staging/}, com os mesmos dados convertidos para Parquet;
	\item o mascaramento de dados pessoais ainda em \texttt{raw/}: CPF e NIS
	são pseudonimizados por HMAC, e nome e data de nascimento são suprimidos,
	em atenção à LGPD \cite{brasil2018lgpd};
	\item camadas bronze, prata e ouro materializadas no \textit{data warehouse}
	por transformações declarativas versionadas, organizadas por domínio, que
	leem o Parquet da \texttt{staging/};
	\item painéis analíticos, que consomem apenas a camada ouro, e um catálogo
	de dados com linhagem;
	\item três ambientes de execução: desenvolvimento local, homologação e
	produção.
\end{itemize}

\begin{figure}[htb]
	\centering
	\resizebox{\textwidth}{!}{%
	\begin{tikzpicture}[
		box/.style={draw, rounded corners, align=center, minimum height=1.2cm,
			text width=3.6cm, font=\footnotesize},
		fonte/.style={box},
		arquivo/.style={box, fill=gray!8},
		bronze/.style={box, fill=orange!18},
		prata/.style={box, fill=gray!15},
		ouro/.style={box, fill=yellow!28},
		orq/.style={box, fill=blue!8},
		gov/.style={box, fill=red!6},
		arr/.style={-{Stealth[length=2mm]}, thick},
		elbl/.style={font=\scriptsize\itshape, fill=white, inner sep=1pt},
		zone/.style={draw, dashed, rounded corners, inner sep=0.3cm},
		lbl/.style={font=\scriptsize\itshape}
	]
		% Fontes
		\node[fonte] (sftp) at (0,6.6)   {SFTP da Caixa\\(FAR, FDS, Rural)};
		\node[fonte] (shp)  at (5.6,6.6) {Planilhas do Novo MCMV\\e histórico da SNH};
		\node[fonte] (apis) at (11.2,6.6) {APIs públicas\\(IBGE, BACEN, CAGED\ldots)};

		% Orquestração
		\node[orq] (airflow) at (5.6,4.2) {Orquestrador\\(Apache Airflow)};

		% Zona de pouso
		\node[arquivo] (raw)     at (0,1.6)  {\textbf{raw/}\\formato original;\\dados pessoais\\mascarados};
		\node[arquivo] (staging) at (0,-1.0) {\textbf{staging/}\\Parquet, com\\colunas de linhagem};

		% Camadas medalhão
		\node[bronze] (bronze) at (5.6,1.6)  {\textbf{Bronze}\\cópia da staging,\\ainda bruta};
		\node[prata]  (prata)  at (5.6,-1.0) {\textbf{Prata}\\tipada e\\unificada por frente};
		\node[ouro]   (ouro)   at (5.6,-3.6) {\textbf{Ouro}\\indicadores\\para consumo};

		% Governança e consumo
		\node[gov] (om)       at (11.2,0.3)  {Catálogo de dados\\(OpenMetadata)};
		\node[box] (superset) at (11.2,-3.6) {Visualização\\(Apache Superset)};

		% Fontes -> orquestrador
		\draw[thick] (sftp.south) |- (5.6,5.35);
		\draw[thick] (apis.south) |- (5.6,5.35);
		\draw[thick] (shp.south) -- (5.6,5.35);
		\draw[arr] (5.6,5.35) -- node[elbl, right] {extração} (airflow.north);

		% Fluxo de dados
		\draw[arr] (airflow.west) -| node[elbl, above, pos=0.25] {grava} (raw.north);
		\draw[arr] (raw) -- node[elbl, right] {converte} (staging);
		\draw[arr] (staging.east) -- ++(1.1,0) |- node[elbl, sloped, above, pos=0.3] {lê o Parquet} (bronze.west);
		\draw[arr] (bronze) -- (prata);
		\draw[arr] (prata) -- (ouro);
		\draw[arr] (ouro) -- node[elbl, above] {leitura} (superset);

		\begin{scope}[on background layer]
			\node[zone, fit=(sftp)(shp)(apis)] (zf) {};
			\node[zone, fit=(raw)(staging)] (zl) {};
			\node[zone, fit=(bronze)(prata)(ouro)] (zw) {};
		\end{scope}

		% Execução e governança
		\draw[arr] (airflow.south) -- node[elbl, right] {executa as transformações} (zw.north);
		\draw[arr, dashed] (om.west) -- node[elbl, above] {cataloga} (om.west -| zw.east);

		\node[lbl, anchor=south west] at (zf.north west) {fontes};
		\node[lbl, anchor=north west, align=left] at (zl.south west) {zona de pouso\\(armazenamento de objetos)};
		\node[lbl, anchor=north west, align=left] at (zw.south west) {\textit{data warehouse}\\(camadas medalhão)};
	\end{tikzpicture}%
	}
	\caption{Arquitetura da plataforma Gov Hub Cidades.}
	\label{fig:govhub}
	\fonte{Elaborado pelos autores, a partir do repositório da plataforma.}
\end{figure}

As ferramentas que compõem essa plataforma, e as que foram adicionadas para o
ML, estão na Tabela~\ref{tab:ferramentas-caso}.

\pendente{Descrever a equipe (papéis, número de pessoas) e o vínculo formal
com o MCID; confirmar com os orientadores o nível de detalhe permitido sobre o
convênio.}

\section{Critérios e procedimentos de avaliação}
\label{sec:avaliacao}

A avaliação combina critérios sobre a arquitetura e sobre a instanciação. A
Tabela~\ref{tab:avaliacao} relaciona cada critério às questões de pesquisa
definidas no Capítulo~\ref{cap:introducao}.

\begin{table}[htb]
	\centering
	\caption{Critérios de avaliação do artefato.}
	\label{tab:avaliacao}
	\small
	\begin{tabularx}{\textwidth}{p{0.9cm}X>{\raggedright\arraybackslash}p{1.8cm}}
		\toprule
		\textbf{ID} & \textbf{Critério e procedimento} & \textbf{Questão} \\
		\midrule
		AV1 & \textbf{Aderência aos princípios de MLOps.} Análise de cada
		princípio P1--P9 de \citeonline{kreuzberger2023}, indicando o
		componente da arquitetura que o atende, ou a justificativa para não
		atendê-lo. & QP, QE1 \\
		AV2 & \textbf{Prontidão para produção.} Aplicação da rubrica
		\textit{ML Test Score} \cite{breck2017} antes e depois da
		implementação. & QP \\
		AV3 & \textbf{Reprodutibilidade.} Reexecução do treino a partir de um
		\textit{commit} e de uma versão de dataset registrados, em ambiente
		limpo, comparando métricas e predições com a execução original. & QE2 \\
		AV4 & \textbf{Ausência de vazamento temporal.} Testes automatizados na
		camada de transformação que verificam que nenhuma feature de uma
		observação no mês $t$ usa
		dado posterior a $t$, e que os cortes de treino, validação e teste
		respeitam a ordem temporal. & QE3 \\
		AV5 & \textbf{Comparação com a regra vigente.} \textit{Recall@k} e
		\textit{precision@k} do modelo e da regra de gargalo no período de
		teste, com $k$ igual à capacidade mensal de fiscalização. & QE3 \\
		\bottomrule
	\end{tabularx}
	\fonte{Elaborado pelos autores.}
\end{table}

O AV5 avalia a instanciação, não a arquitetura: um modelo que não supere a
regra vigente não invalida a arquitetura, desde que a comparação tenha sido
feita de forma correta e reprodutível. Esse resultado negativo seria, ele
próprio, uma informação útil para o órgão.

\pendente{Decidir com os orientadores se haverá uma avaliação qualitativa
adicional com a equipe técnica (por exemplo, questionário ou grupo focal sobre
a facilidade de adicionar um segundo modelo à arquitetura).}

\section{Ferramentas}
\label{sec:ferramentas}

A arquitetura é definida por papéis, independentes de ferramenta
(Seção~\ref{sec:arquitetura}). As ferramentas que cumprem cada papel no caso
estudado são definidas uma única vez, na Tabela~\ref{tab:ferramentas-caso},
junto com a instanciação. A mesma tabela inclui as ferramentas de apoio à
pesquisa: Jupyter, para a análise exploratória, e Git com GitHub Actions, para
versionamento e integração contínua.
